\section{Positive energy and the joint momentum measure}
\label{sec:spectral-measure}
The Fourier support of the vacuum distributions and the spectrum of the
translation operators are different assertions. We now prove the latter.
The essential analytic fact is elementary: a positive finite measure whose
Fourier transform extends boundedly to the upper half-plane cannot charge
negative frequencies. Positivity is what makes a Gaussian probe detect
such a charge without cancellation.

We use the spectral theorem for strongly continuous unitary representations
of locally compact abelian groups, in the form of the SNAG theorem
\cite[Theorem D.3.1, p.397]{BHV}. Applied to a strongly continuous unitary
representation $T$ of $\mathbb R^4$, it supplies a unique regular
projection-valued measure $P$ on $\mathbb R^4$ with
\begin{equation}\label{eq:snag-convention}
 T(a)=\int e^{ip\cdot a}\,dP(p),\qquad P(\mathbb R^4)=I.
\end{equation}
Here we identify the character indexed by $p$ with $a\mapsto e^{ip\cdot a}$;
the pairing is Euclidean, as in Section~\ref{sec:fourier-support}.
The hypotheses hold for the translation representation constructed in
Proposition~\ref{prop:word-motion}, and for its Euclidean counterpart in
Section~\ref{sec:complex-heat}. In particular, continuity is strong operator
continuity, not operator-norm continuity. For each vector $v$,
$\mu_v(B)=\langle v,P(B)v\rangle$ is a positive regular finite measure,
of mass $\|v\|^2$, with characteristic function $\langle v,T(a)v\rangle$.

\begin{lemma}[A bounded upper-half-plane extension]\label{lem:paley-positive}
Let $\mu$ be a positive finite Borel measure on $\mathbb R^d$, and let
$\ell:\mathbb R^d\to\mathbb R$ be continuous. Suppose a function $F$ is
continuous on $\operatorname{Im}z\geq0$, holomorphic on
$\operatorname{Im}z>0$, bounded there by $C$, and satisfies
\[
 F(t)=\int e^{it\ell(p)}\,d\mu(p)\qquad(t\in\mathbb R).
\]
Then $\operatorname{supp}\mu\subseteq\{p:\ell(p)\geq0\}$.
\end{lemma}
\begin{proof}
Fix $\varepsilon>0$, $b\in\mathbb R$ and $s\geq0$. Multiply $F(z)$ by
$\exp(-ibz-\varepsilon z^2)$ and integrate around the rectangle with
corners $-R,R,R+is,-R+is$. Cauchy's theorem applies: the function is
continuous on the rectangle and holomorphic in its interior. One can
also apply the theorem above height $\delta>0$ and let $\delta\downarrow0$;
continuity on each compact rectangle justifies that limit.
On either vertical edge the magnitude is at most
$C\exp(-\varepsilon R^2+|b|s+\varepsilon s^2)$.
The vertical integrals tend to zero. The horizontal integrals converge
absolutely, since at height $u$ the common majorant is
$C\exp(-\varepsilon t^2+bu+\varepsilon u^2)$.
Consequently the integrals along heights zero and $s$ agree, and
\begin{equation}\label{eq:gaussian-shift-bound}
 \left|\int_{\mathbb R}F(t)e^{-ibt-\varepsilon t^2}\,dt\right|
 \leq C\sqrt{\frac\pi\varepsilon}\,
                e^{bs+\varepsilon s^2}.
\end{equation}

Fubini applies to the boundary integral because $\mu$ is finite and
$e^{-\varepsilon t^2}$ is integrable. The Gaussian transform gives
\[
 \int_{\mathbb R}e^{it(\ell(p)-b)-\varepsilon t^2}\,dt
 =\sqrt{\frac\pi\varepsilon}\,
       e^{-(\ell(p)-b)^2/(4\varepsilon)}.
\]
For completeness, if $G(a)=\int e^{iat-\varepsilon t^2}\,dt$,
differentiation under the integral and integration by parts give
$G'(a)=-aG(a)/(2\varepsilon)$. The value
$G(0)=\sqrt{\pi/\varepsilon}$ follows by squaring the Gaussian integral
and using polar coordinates. Solving this differential equation proves
the formula with its stated normalization.

Suppose $p_0\in\operatorname{supp}\mu$ and $b=\ell(p_0)<0$.
In \eqref{eq:gaussian-shift-bound} choose $s=-b/(2\varepsilon)$ and
cancel the positive factor $\sqrt{\pi/\varepsilon}$. Positivity of the
resulting integral yields
\[
 \int e^{-(\ell(p)-b)^2/(4\varepsilon)}\,d\mu(p)
 \leq C e^{-b^2/(4\varepsilon)}.
\]
The open neighborhood
$A=\{p:(\ell(p)-b)^2<b^2/4\}$ of $p_0$ has mass $m>0$.
Its contribution to the left side is at least
$m e^{-b^2/(16\varepsilon)}$. Hence
$m\leq C e^{-3b^2/(16\varepsilon)}$ for every $\varepsilon>0$.
Letting $\varepsilon\downarrow0$ is a contradiction.
\end{proof}

Apply the lemma first on the Euclidean Hilbert space. Write
$T_E(t,\mathbf x)=V(t)U_E(\mathbf x)$ for its joint translations,
where $V$ is the unitary time group of Proposition~\ref{prop:complex-heat}.
For $v\in\mathcal H_E$ define
\[
 F_v(t+is)=\langle v,V(t)E_s v\rangle\qquad(s\geq0).
\]
This function is continuous on the closed upper half-plane by strong
continuity and the uniform operator bounds. On its interior it equals
$\langle v,E(-iz)v\rangle$, so is holomorphic and bounded by $\|v\|^2$.
Its boundary is the characteristic function of the time coordinate of
the joint measure $\mu^E_v$. The lemma gives
\begin{equation}\label{eq:euclidean-positive-energy}
 \operatorname{supp}\mu^E_v\subseteq\{p:p^0\geq0\}.
\end{equation}
Spatial covariance is not used in this step.

\begin{theorem}[Forward spectrum and the unique vacuum]\label{thm:field-spectrum}
Suppose the word construction has a strongly continuous Poincar\'e
representation $U$ with proper orthochronous Lorentz part, and an isometric
translation intertwiner $\mathcal V:\mathcal H\to\mathcal H_E$ as in
Proposition~\ref{prop:word-motion}. Then its joint momentum measure satisfies
\[
 P(\mathbb R^4\setminus\overline V_+)=0,
 \qquad \overline V_+=\{p:p^0\geq|\mathbf p|\}.
\]
If the heat-fixed subspace of $\mathcal H_E$ is
$\mathbb C\Omega_E$, then the vectors fixed by all translations on
$\mathcal H$ are exactly $\mathbb C\Omega$.
\end{theorem}
\begin{proof}
The intertwining identity and isometry give
\[
 \langle v,U(a)v\rangle
 =\langle\mathcal Vv,T_E(a)\mathcal Vv\rangle.
\]
Thus $\mu_v$ and $\mu^E_{\mathcal Vv}$ have equal characteristic functions and
are equal measures. Here is the measure-uniqueness step explicitly.
For any Schwartz function $f$, Fourier inversion and Fubini express its
integral against a finite measure in terms of that characteristic
function; Fubini follows from integrability of $\widehat f$.
The two measures therefore agree on all smooth compactly supported tests.
Continuous compactly supported tests are uniformly approximated, in a
fixed compact neighborhood of their supports, by smooth mollifications.
They too have equal integrals. Two finite positive Borel measures on
$\mathbb R^4$ with equal integrals on $C_c(\mathbb R^4)$ are equal. Indeed,
for compact $K$ the functions $\phi_k(p)=\max(0,1-k\operatorname{dist}(p,K))$
lie in $C_c(\mathbb R^4)$ and decrease pointwise to the indicator of $K$, so
dominated convergence gives equal masses on every compact set; compact sets
form a $\pi$-system generating the Borel $\sigma$-algebra, and the total
masses agree by monotone convergence along balls, so Dynkin's
$\pi$--$\lambda$ theorem gives equality of the measures. In particular every
$\mu_v$ obeys \eqref{eq:euclidean-positive-energy}.

For a Lorentz matrix $\Lambda$, covariance of translations implies
\[
 \mu_{U(\Lambda)v}=(\Lambda^{-T})_*\mu_v.
\]
Indeed the characteristic function of the left side at $a$ equals
$\langle v,U(\Lambda^{-1}a)v\rangle$, which is that of the right side.
The preceding uniqueness argument applies again. A homeomorphism carries
the support of a measure to the support of its pushforward: every
neighborhood of the image has a preimage neighborhood of positive measure,
and the converse follows by its inverse. It follows that, for any
$p\in\operatorname{supp}\mu_v$,
\[
 (\Lambda^{-T}p)^0\geq0\quad\hbox{for every proper orthochronous }\Lambda.
\]
The dual matrices $\Lambda^{-T}=\eta\Lambda\eta$, with
$\eta=\operatorname{diag}(1,-1,-1,-1)$, range over the same group.
The identity gives $p^0\geq0$. If $|\mathbf p|>p^0$, take a boost in the
direction $\mathbf p/|\mathbf p|$ with positive rapidity $r$ such that
$\tanh r>p^0/|\mathbf p|$. Its time coordinate is
$p^0\cosh r-|\mathbf p|\sinh r<0$, a contradiction. Thus every
$\mu_v$ is supported in $\overline V_+$.
The complement of the support of a finite Borel measure on $\mathbb R^4$
has measure zero: it is a union of a countable collection of zero-measure
basis neighborhoods. Hence $\mu_v(\overline V_+^c)=0$ for every $v$.
Since $P(B)$ is an orthogonal projection,
$\|P(B)v\|^2=\mu_v(B)$, which proves the asserted operator identity.

Finally let $v$ be translation invariant. Its image $\mathcal Vv$ is fixed by all
Euclidean real-time translations. The fixed-vector equivalence proved in
Section~\ref{sec:complex-heat} makes it heat fixed. Therefore
$\mathcal Vv=c\Omega_E=\mathcal V(c\Omega)$; injectivity of the isometry gives
$v=c\Omega$. The converse follows from vacuum invariance. This proves
uniqueness under the stated heat hypothesis, and makes no spectral-gap
assertion.
\end{proof}
